% ============================================================================
% CHAPITRE 4 — SECTION 4.3 : DISCUSSIONS
% Fichier : partie4/discussions.tex
% Version : 2.2 — alignée sur la formalisation Lean4 révisée
% ============================================================================

\section{DISCUSSIONS}
\label{sec:discussions}
% ============================================================================
\begin{sloppypar}
La
section~\ref{sec:resultats_theoriques} a mis le cadre CLO-CRO-PQ-SOTP en
regard des estimations analytiques disponibles pour cinq syst\`emes d\'eploy\'es
(trois classiques et deux signatures post-quantiques), dans une calibration
illustrative. Elle a dress\'e le bilan d'un corpus th\'eorique articul\'e selon la
progression F$\to$B$\to$E$\to$S$\to$C$\to$X.
\'Etablir ces faits ne suffit toutefois pas \`a en fixer la port\'ee. La pr\'esente section
examine ce qu'ils signifient au-del\`a de la cryptographie institutionnelle consid\'er\'ee en vase
clos, c'est-\`a-dire pour l'avancement de la cryptographie comme discipline, pour les
praticiens qui con\c{c}oivent des syst\`emes et pour les probl\`emes ouverts
qui restent non r\'esolus.

L'analyse proc\`ede en onze temps. Elle examine d'abord ce que la
formalisation a r\'ev\'el\'e sur le Trilemme~CRO
(sous-section~\ref{subsec:ce_que_formalisation_revele}), puis positionne nos
r\'esultats parmi les impossibilit\'es fondamentales en pr\'esentant la
classification \'epist\'emologique des contributions
(sous-section~\ref{subsec:positionnement_impossibilites}). Elle d\'etaille ensuite l'apport
sp\'ecifique de l'extension post-quantique
(sous-section~\ref{subsec:apport_extension_quantique}), d\'egage les quatre sch\'emas
structurels du compromis CLO (sous-section~\ref{subsec:quatre_patterns}) et mesure la sensibilit\'e de
l'opposabilit\'e \`a la capacit\'e d'extraction
(sous-section~\ref{subsec:analyse_sensibilite}). Les implications pour
la conception et la r\'eglementation (sous-section~\ref{subsec:implications_conception})
pr\'ec\`edent la situation de la contribution
vis-\`a-vis des travaux existants (sous-section~\ref{subsec:positionnement_travaux}) ainsi qu'un bilan
des limitations et probl\`emes ouverts
(sous-section~\ref{subsec:limitations_problemes_ouverts}). Les trois derni\`eres sous-sections montrent comment la
famille SOTP transforme l'impossibilit\'e en guide
(sous-section~\ref{subsec:portee_constructive}), tirent la le\c{c}on
m\'ethodologique de la certification
(sous-section~\ref{subsec:verification_corrige_lecon}) et clarifient le
statut \'epist\'emique du programme exp\'erimental
(sous-section~\ref{subsec:statut_protocoles_experimentaux}).
\end{sloppypar}

\begin{boiteArbitre}
L'arbitre dispose maintenant des faits, \`a savoir des bornes prouv\'ees sous hypoth\`eses,
des constructions dont les propri\'et\'es cryptographiques sont \'etablies mais
dont l'opposabilit\'e reste heuristique ainsi qu'une formalisation qui a mis au jour
des lacunes du raisonnement initial. Ce
chapitre lui demande ce que tout cela \emph{signifie} pour la discipline
naissante, pour les praticiens qui choisissent des protocoles et pour les
questions qui restent ouvertes. Une impossibilit\'e prouv\'ee constitue \`a ses yeux une
forme nette de connaissance, puisqu'elle lui apprend ce qu'il ne pourra jamais obtenir
et pour quelle raison.
\end{boiteArbitre}

% -----------------------------------------------------------------------
\subsection{Ce que la formalisation a r\'ev\'el\'e}
\label{subsec:ce_que_formalisation_revele}
% -----------------------------------------------------------------------
\begin{sloppypar}
La v\'erification m\'ecanique du Trilemme~CRO dans Lean4 a produit,
au-del\`a de la certification du c\oe ur alg\'ebrique, cinq r\'ev\'elations
\'epist\'emologiques que le raisonnement informel n'avait pas permis
d'identifier.

\subsubsection{Hypoth\`ese explicite des propri\'et\'es d'instance}
\label{subsubsec:hypothese_quasi_optimalite}

La preuve originale~\cite{MINKA2026CRO} utilise les propri\'et\'es
informationnelles (Fano, Pinsker) sans les expliciter comme conditions
v\'erifiables du th\'eor\`eme. La formalisation Lean4 les rend
explicites en les int\'egrant comme champs de la structure
\texttt{ProtocolInstance}~:
\begin{equation}
\texttt{pinsker} : \TV^2 \leq c_0 \cdot \KL, \qquad
\texttt{fano} : I \geq \Relrec \cdot \lambda - 1.
\label{eq:hypothese_quasi_optimalite}
\end{equation}
Cette explicitation s'aligne avec la rationalit\'e
born\'ee~\cite{Simon82,Miller56,Ericsson95}, puisque les institutions ne peuvent
extraire qu'une information limit\'ee d'une preuve cryptographique, ce qui
force les protocoles \`a op\'erer dans des r\'egimes de faible information
relativement \`a $\lambda$. La formalisation Lean4 transforme des
intuitions diffuses en conditions v\'erifiables.

\subsubsection{Structure en trois \'etapes de la preuve formalis\'ee}
\label{subsubsec:non_derivabilite_borne}

La formalisation \'etablit la borne
$\gamma^{-1}(\Delta/\lambda + \kappa(\Delta+1)/\lambda^2 + \kappa^2/\lambda^3)$
en trois \'etapes principales. (1)~Fano borne $\Relrec$~; (2)~l'extraction
born\'ee et le verdict efficace bornent $\Opp$~; (3)~on forme le produit
puis on applique le r\'egime de confidentialit\'e
$I(M;\Sigma|C) \leq \kappa/\lambda$. Cette structure, plus simple que
l'optimisation sur $[0,\lambda]$, pr\'eserve la borne dominante
$\Delta/\lambda$ et la propri\'et\'e qualitative centrale du trilemme.
La formalisation compl\`ete de l'optimisation analytique sur
$[0,\lambda]$ n\'ecessiterait des outils d'analyse r\'eelle non encore
disponibles dans Mathlib~\cite{MathlibAnalysis,Tsybakov2009,Scerri14}.

\subsubsection{N\'egligeabilit\'e~: une preuve par limite}
\label{subsubsec:necessite_analyse_cas}

L'\'enonc\'e informel \og $(1/2)^\lambda$ est n\'egligeable\fg{} est tenu
pour \'evident dans la litt\'erature
cryptographique~\cite{Goldreich01,KatzLindell14}. Dans la formalisation, la
n\'egligeabilit\'e est d\'efinie par un filtre (\og pour tout $c$, \`a partir d'un
certain rang, $\varepsilon(\lambda)\,\lambda^c \leq 1$\fg)~; la preuve passe par
la limite $\lambda^c/2^\lambda \to 0$ fournie par Mathlib. Elle n'est pas
constructive, puisqu'elle ne produit pas de seuil explicite. La formulation initiale
de ce paragraphe d\'ecrivait une analyse de cas ($c \le 4$ / $c > 4$) et un seuil
$N(c) = \lceil 2c \rceil + 4$~; la pr\'esente formulation suit le code compil\'e. Le ph\'enom\`ene relev\'e par Scerri~\cite{Scerri14} demeure, car
formaliser un fait asymptotique \og \'evident\fg{} oblige \`a choisir une
d\'efinition pr\'ecise.

\subsubsection{Transparence des constantes implicites}
\label{subsubsec:transparence_constantes}

La notation $O(\lambda^{-2})$ dissimule une d\'ependance sur
$1/\gamma$ o\`u $\gamma = \alpha H(C) - 1$. Notre borne explicite
$\gamma^{-1}(\Delta/\lambda + \kappa(\Delta+1)/\lambda^2 +
\kappa^2/\lambda^3)$ expose cette d\'ependance, puisque la constante
d\'epend \`a la fois de $\Delta$ (capacit\'e d'extraction
institutionnelle) et de $\gamma$ (granularit\'e des verdicts
efficace). Pour la cryptographie
institutionnelle, o\`u les param\`etres doivent \^etre empiriquement
mesurables, cette transparence constitue une
exigence op\'erationnelle.

\subsubsection{Les param\`etres institutionnels sont des constantes du th\'eor\`eme}
\label{subsubsec:parametres_constantes}

La formalisation \'etablit que $\Delta$, $\alpha$, $\beta$, $H(C)$ ont le
statut de constantes
du th\'eor\`eme, distinct de celui de param\`etres libres, de sorte que leur mesure empirique devient un
pr\'erequis pour son application. Un syst\`eme d\'eployant un protocole
sans mesurer pr\'ealablement sa capacit\'e d'extraction $\Delta$ ne peut
pas \'evaluer si la borne th\'eorique est satisfaite~\cite{Grossklags10,Herley09}.

\subsubsection{Coh\'erence et port\'ee du th\'eor\`eme principal}
\label{subsubsec:consequences_empiriques_borne}

La formalisation Lean4 produit deux t\'emoins qui d\'emontrent la
non-vacuit\'e et la n\'ecessit\'e des hypoth\`eses. Le \textbf{t\'emoin
de coh\'erence} est une instance qui v\'erifie toutes les hypoth\`eses
avec $\Conf = \Opp = 1$, $\Rel = 1/\lambda$, $I = 0$~; pour cette
instance, $\Conf \cdot \Rel \cdot \Opp = 1/\lambda$, ce qui \'egale
exactement la borne du th\'eor\`eme. Il s'agit d'un point du mod\`ele abstrait,
qu'aucun protocole ne r\'ealise (avec $I = 0$, $\Relrec = 2^{-\lambda}$~;
Remarque~\ref{rem:portee_temoin}). Le \textbf{t\'emoin de port\'ee}
est une instance qui v\'erifie les Hypoth\`eses 1, 2$'$ et 3 avec
$\Delta = 0$ mais s'\'ecarte du r\'egime de confidentialit\'e
($I = \lambda/2 - 1$)~; elle pr\'esente $\Conf \cdot \Rel \cdot \Opp = 1/4$
constant. Ces deux t\'emoins \'etablissent que les hypoth\`eses sont
satisfiables et que le r\'egime est n\'ecessaire \`a la borne.
\end{sloppypar}

\begin{boiteArbitre}
L'arbitre s'arr\^ete sur ce que la machine a trouv\'e, soit cinq points que le
raisonnement \`a la main n'avait pas vus. Des propri\'et\'es
informationnelles y deviennent des conditions v\'erifiables, une
structure de preuve pr\'eserve la borne dominante et un
fait \og \'evident \fg{} exige une d\'efinition pr\'ecise. S'y ajoutent des constantes que
la notation $O(\cdot)$ cachait ainsi que des param\`etres institutionnels qui
se r\'ev\`elent \^etre des constantes du th\'eor\`eme, distinctes de variables
libres.

Il note que la v\'erification formelle a fait davantage que confirmer le
r\'esultat, puisqu'elle l'a approfondi, ce qu'il tient pour un gain r\'eel.
\end{boiteArbitre}

% -----------------------------------------------------------------------
\subsection{Positionnement parmi les impossibilit\'es fondamentales}
\label{subsec:positionnement_impossibilites}
% -----------------------------------------------------------------------
\begin{sloppypar}

\subsubsection{La cryptographie institutionnelle comme domaine propre}
\label{subsubsec:crypto_institutionnelle_domaine}

La m\'etrique $\Opp = I(V;\Sigma|C)/H(V|C)$ est absente de tout cadre
cryptographique ant\'erieur. Le calcul multi-parties
s\'ecuris\'e~\cite{Yao82,GMW87,Goldreich1987,Evans2018} optimise le
nombre de rounds et la confidentialit\'e du calcul distribu\'e~; le
chiffrement homomorphe~\cite{Gentry09,Vaikuntanathan2012,Brakerski2014}
optimise la profondeur de circuits sur donn\'ees chiffr\'ees~; les
syst\`emes zero-knowledge~\cite{GMR85,Groth16,BCTV14} minimisent la
fuite d'information tout en prouvant l'appartenance. Aucune de ces
sp\'ecialisations ne pose la question de savoir \emph{dans quelle mesure un
acteur institutionnel born\'e cognitivement peut extraire une
information utile de cette preuve pour rendre un verdict
inform\'e}~\cite{Simon82,Miller56,Kahneman11}.

Les travaux fondateurs de Goldwasser-Micali-Rackoff~\cite{GMR85} ont
axiomatis\'e la puissance computationnelle de
l'adversaire~\cite{Goldreich01,Bellare97}~; nous axiomatisons ici la
capacit\'e interpr\'etative de l'interpr\`ete. Il s'agit du m\^eme objet, un acteur
born\'e, consid\'er\'e selon une sym\'etrie invers\'ee.

\begin{table}[htbp]
\centering
\caption{Positionnement de la cryptographie institutionnelle parmi les
sp\'ecialisations cryptographiques}
\label{tab:comparaison_domaines_crypto}
\begin{tabular}{@{}p{3.5cm}p{3.5cm}p{4.5cm}@{}}
\toprule
\textbf{Domaine} & \textbf{Optimise pour} & \textbf{M\'etrique centrale} \\
\midrule
Cryptographie classique   & R\'esistance adversaires
  & Avantage distingueur~\cite{Goldreich01} \\
MPC                       & Calcul distribu\'e priv\'e
  & Rounds, parties~\cite{Yao82,Evans2018} \\
FHE                       & Calcul sur donn\'ees chiffr\'ees
  & Profondeur circuits~\cite{Gentry09,Brakerski2014} \\
ZKP                       & Preuves sans r\'ev\'elation
  & Taille preuve, compl\'etude~\cite{GMR85,Groth16} \\
\textbf{Crypto institutionnelle}   & \textbf{Interpr\'etabilit\'e l\'egale}
  & \textbf{Opposabilit\'e $\Opp$}~\cite{MINKA2026CLO,MINKA2026CRO} \\
\textbf{Crypto PQ institutionnelle} & \textbf{R\'esilience quantique}
  & \textbf{$\Delta_q$, $\eta_q$}~\cite{MINKA2026PQCRO} \\
\bottomrule
\end{tabular}
\end{table}

\subsubsection{Le Trilemme CRO parmi les impossibilit\'es fondamentales}
\label{subsubsec:trilemme_parmi_impossibilites}

Le Trilemme CRO conditionnel joue dans la cryptographie institutionnelle
le m\^eme r\^ole structurant que les th\'eor\`emes d'impossibilit\'e
dans d'autres branches de l'informatique
th\'eorique~\cite{Abadi2012,Attiya2004}. Le th\'eor\`eme
FLP~\cite{FLP85} \'etablit qu'aucun protocole d\'eterministe
ne garantit le consensus en pr\'esence d'une seule d\'efaillance
asynchrone~; le th\'eor\`eme CAP~\cite{Brewer00,GilbertLynch02} \'etablit
qu'un syst\`eme distribu\'e ne peut simultan\'ement garantir coh\'erence
stricte, disponibilit\'e totale et tol\'erance aux partitions~; le
Trilemme CRO \'etablit, sous ses hypoth\`eses, que le produit $\Conf\cdot\Relrec\cdot\Opp$ est
born\'e par
$\gamma^{-1}(\Delta/\lambda + \kappa(\Delta+1)/\lambda^2 + \kappa^2/\lambda^3)$
(classique) ou par son analogue quantique
$\gamma_q^{-1}(\Delta_q/\lambda + \kappa(\Delta_q+1)/\lambda^2 + \kappa^2/\lambda^3)$
(\cite{MINKA2026CRO} et section~\ref{sec:cro_pq}).

\begin{table}[htbp]
\centering
\caption{Comparaison des th\'eor\`emes d'impossibilit\'e fondamentaux}
\label{tab:comparaison_impossibilites}
\begin{tabular}{@{}p{2.8cm}p{4cm}p{4.5cm}@{}}
\toprule
\textbf{Th\'eor\`eme} & \textbf{Contrainte fondamentale}
  & \textbf{Hypoth\`eses} \\
\midrule
FLP & Consensus impossible (asynchrone + 1 d\'efaillance)
  & Mod\`ele asynchrone~\cite{FLP85,Attiya2004} \\
CAP & C\,+\,A\,+\,P impossibles simultan\'ement
  & Syst\`eme distribu\'e~\cite{Brewer00,GilbertLynch02} \\
Blockchain Trilemma & D\'ecentralisation + S\'ecurit\'e + Scalabilit\'e
  & Blockchain~\cite{Buterin2014,Abadi2012} \\
\textbf{CRO classique} & $\Conf\cdot\Relrec\cdot\Opp \le \gamma^{-1}(\Delta/\lambda + \kappa(\Delta+1)/\lambda^2 + \kappa^2/\lambda^3)$
  & \textbf{$\Delta,\alpha,\beta$ institutionnels}~\cite{MINKA2026CRO} \\
\textbf{CRO post-quantique} &
  $\Conf^q\!\cdot\!\Relrec^q\!\cdot\!\Opp^q \le
   \gamma_q^{-1}(\Delta_q/\lambda + \kappa(\Delta_q+1)/\lambda^2 + \kappa^2/\lambda^3)$
  & \textbf{$\Delta_q$, $\eta_q$ quantiques}~\cite{MINKA2026PQCRO} \\
\bottomrule
\end{tabular}
\end{table}

Une distinction s'impose toutefois, car FLP et CAP ont \'et\'e prouv\'es dans des
mod\`eles de calcul formels bien
\'etablis alors que le Trilemme CRO est conditionnel
aux hypoth\`eses sur $\Delta$, $\alpha$, $\beta$, $\Delta_q$,
$\eta_q$, param\`etres dont les valeurs n'ont pas encore \'et\'e
valid\'ees par des protocoles exp\'erimentaux contr\^ol\'es sur des
acteurs institutionnels r\'eels. Le Trilemme
CRO est un th\'eor\`eme correct sous ses hypoth\`eses~; sa pertinence
empirique reste \`a \'etablir.

Sous cette r\'eserve, le d\'eplacement de perspective est net. Que
$\Delta \in [35, 73]$~bits s'av\`ere plus restrictif que
$\lambda \in [128, 256]$~bits signifie que la contrainte dominante sur
les syst\`emes cryptographiques institutionnels tient \`a la capacit\'e
interpr\'etative de l'institution, probl\`eme largement ignor\'e
jusqu'ici. La puissance de calcul de l'adversaire, que la cryptographie sait
traiter, passe au second plan.

\subsubsection{Classification \'epist\'emologique des r\'esultats}
\label{subsubsec:classification_epistemologique}

Les r\'esultats de cette th\`ese, qui ne sont pas de m\^eme nature, doivent
\^etre distingu\'es selon leur niveau de
certitude.

\begin{table}[htbp]
\centering
\caption{Hi\'erarchie de certitude des r\'esultats de la th\`ese}
\label{tab:classification_epistemologique}
\begin{tabular}{@{}p{3.8cm}p{3.2cm}p{4.5cm}@{}}
\toprule
\textbf{R\'esultat} & \textbf{Niveau de certitude} & \textbf{Statut} \\
\midrule
R\'esultats CLO (bornes, goulot, composition) & Prouv\'es (formulations pr\'ecis\'ees)
  & Preuves papier~(section~\ref{sec:clo}) \\
Trilemme classique (forme raffin\'ee) & Sous hypoth\`eses institutionnelles
  & Noyau v\'erifi\'e Lean4~(\cite{MINKA2026CRO} et section~\ref{sec:lean4}) \\
Trilemme PQ & Sous hypoth\`eses quantiques
  & Preuve papier~(\cite{MINKA2026PQCRO} et section~\ref{sec:cro_pq}) \\
Propri\'et\'es SOTP ($\Conf$, $\Relver$) & Sous PRG, DLog, q-PKE
  & R\'eductions~(annexe~\ref{chap:annexe_gamebased}) \\
Opposabilit\'e SOTP, optimalit\'e & Heuristique / conjecture
  & Non prouv\'ees~(Conjecture~\ref{conj:optimalite_universelle}) \\
Formalisation Lean4 & M\'ecaniquement certifi\'ee
  & D\'erivation v\'erifi\'ee, pr\'emisses pos\'ees~(section~\ref{sec:lean4}) \\
\midrule
\multicolumn{3}{l}{\textit{Falsifiables exp\'erimentalement~:}} \\
Valeurs de $\Delta$, $\Delta_q$, $\eta_q$ & Empirique
  & \`A valider par \'etudes contr\^ol\'ees~(chapitre~\ref{chap:experimentation}) \\
\bottomrule
\end{tabular}
\end{table}

Distinguer explicitement ces niveaux est une contribution
m\'ethodologique propre~\cite{Nosek2015,IOC2019}, puisqu'elle prot\`ege contre
la sur\'evaluation des r\'esultats conditionnels et ouvre pr\'ecis\'ement
le programme exp\'erimental futur. La fronti\`ere entre \og th\'eor\`eme
correct sous hypoth\`eses \fg{} et \og fait empirique \'etabli \fg{} doit
\^etre maintenue avec rigueur~\cite{Lakatos1978}.
\end{sloppypar}

% -----------------------------------------------------------------------
\subsection{L'apport de l'extension quantique}
\label{subsec:apport_extension_quantique}
% -----------------------------------------------------------------------
\begin{sloppypar}
L'extension quantique (section~\ref{sec:cro_pq}) introduit deux
innovations th\'eoriques sans analogue dans le r\'egime classique.

\subsubsection{P\'enalit\'e s\'emantique quantique}
\label{subsubsec:penalite_semantique}

Le terme $\eta_q \cdot H(C)/\lambda$ repr\'esente une contrainte
fondamentalement nouvelle, puisqu'il chiffre la part de la preuve qui, sous
bruit de canal quantique, ne peut plus \^etre interpr\'et\'ee. Il
n'intervient pas dans la borne explicite. Il est aussi nul pour toute preuve
qui est une cha\^{\i}ne de bits classique, y compris les signatures ML-DSA et
SLH-DSA. Il ne devient pertinent que dans le sc\'enario, hypoth\'etique \`a ce
jour, d'un support ou d'un v\'erificateur quantique, o\`u il fournit une
pr\'ediction mesurable.
Cette p\'enalit\'e \'emerge de la th\'eorie de l'information quantique
(borne de Holevo~\cite{Holevo1973}, lemme de la mesure
douce~\cite{Winter1999,Ozawa2000}).

\subsubsection{D\'ecomposition de la capacit\'e quantique}
\label{subsubsec:decomposition_capacite_quantique}

La formule $\Delta_q = \Delta + \delta_{\text{enh}} - \delta_{\text{degr}}$
d\'ecompose explicitement les effets quantiques sur la capacit\'e
institutionnelle~:
\begin{equation}
\Delta_q = \Delta + \delta_{\text{enh}} - \delta_{\text{degr}}
\label{eq:decomposition_delta_q_discussion}
\end{equation}
\begin{itemize}
  \item $\delta_{\text{enh}}$ d\'esigne l'am\'elioration par outils quantiques
    (recherche Grover~\cite{Grover96}, apprentissage machine
    quantique~\cite{Biamonte2017})~;
  \item $\delta_{\text{degr}}$ d\'esigne la d\'egradation par bruit
    quantique~\cite{Preskill1998,Devoret2013} et inadaptation
    institutionnelle.
\end{itemize}
Cette d\'ecomposition permet de mod\'eliser l'\'evolution temporelle de
$\Delta_q$ \`a mesure que les institutions s'adaptent aux syst\`emes
post-quantiques et que les technologies quantiques
m\^urissent. Les valeurs retenues, $\Delta_q \in [25, 67]$~bits, sont
des hypoth\`eses de travail que le protocole de la
section~\ref{subsec:protocole_delta_q} permettra de mesurer.
\end{sloppypar}

% -----------------------------------------------------------------------
\subsection{Quatre sch\'emas fondamentaux du compromis CLO}
\label{subsec:quatre_patterns}
% -----------------------------------------------------------------------
\begin{sloppypar}

Les valeurs d\'eriv\'ees au chapitre~\ref{chap:experimentation}
(tableau~\ref{tab:composantes_opposabilite}) r\'ev\`elent quatre sch\'emas
structurels r\'ecurrents.

\medskip
\noindent\textbf{Sch\'ema 1. Tension vie priv\'ee--opposabilit\'e.}
Les syst\`emes maximisant la confidentialit\'e (Helios $\Conf = 0{,}92$,
zk-SNARK $\Conf = 0{,}89$) atteignent une opposabilit\'e mod\'er\'ee
($\Opp \approx 0{,}50$--$0{,}67$), tandis que le syst\`eme \`a haute
opposabilit\'e (ECDSA $\Opp = 0{,}90$) sacrifie la confidentialit\'e
($\Conf = 0{,}31$). Ce sch\'ema mat\'erialise la tension d\'ecrite par le
paradoxe de l'authenticit\'e invisible~\cite{MINKA2026CLO}.

\noindent\textit{Formalisation.} Ce sch\'ema est captur\'e par le Lemme
Conf-MI (sous-section~\ref{subsec:proprietes_fondamentales}), selon lequel plus
$I(M;\Sigma|C)$ est grand, plus $\Conf$ est faible. ECDSA op\`ere \`a une
extr\'emit\'e de ce compromis ($I$ grand, $\Conf$ faible). Helios est \`a
l'autre extr\'emit\'e pour la confidentialit\'e, mais sa fiabilit\'e
$\Rel = 0{,}95$ est une fiabilit\'e de \emph{v\'erification}, car sous $\Relrec$ Fano
imposerait $I(M;\Sigma|C) \geq 0{,}95\lambda - 1 \approx 121$~bits, incompatible
avec une confidentialit\'e \'elev\'ee. Le sch\'ema illustre donc la tension, sans
en \^etre une instance.

\medskip
\noindent\textbf{Sch\'ema 2. D\'ependance au contexte institutionnel.}
L'opposabilit\'e varie significativement avec l'expertise de
l'institution. Les r\'egulateurs sp\'ecialis\'es
($H(\mathcal{C}) = 11$ bits) extraient plus d'information que les
tribunaux g\'en\'eraux ($H(\mathcal{C}) = 8$ bits), tandis que les
tribunaux dot\'es d'une expertise PKI ($H(\mathcal{C}) = 20$ bits)
atteignent une interpr\'etabilit\'e substantiellement plus \'elev\'ee.
$H(\mathcal{C})$ est donc un param\`etre de conception critique, sur
lequel le concepteur peut agir. Les valeurs cit\'ees ici sont celles du contexte
d'esp\`ece~; elles ne se confondent pas avec les entropies de domaine du
tableau~\ref{tab:estimations_structurelles_HC}, distinction \'etablie
\`a la section~\ref{sec:methodologie_calcul}.

\noindent\textit{Implication pratique.} D\'eployer un zk-SNARK devant un
tribunal g\'en\'eral ($H(\mathcal{C}) \approx 8$ bits) au lieu d'un
r\'egulateur sp\'ecialis\'e ($H(\mathcal{C}) \approx 11$ bits)
r\'eduirait $I(V;\Sigma|C)$ de $2$ bits \`a environ $1$ bit ainsi que
l'opposabilit\'e de $0{,}67$ \`a environ $0{,}33$.

\medskip
\noindent\textbf{Sch\'ema 3. Plafond d'extraction des syst\`emes
z\'ero-connaissance.} M\^eme avec une entropie contextuelle \'elev\'ee,
les syst\`emes cryptographiquement opaques limitent fondamentalement
l'extraction institutionnelle \`a un ou deux bits. Ce plafond manifeste
l'hypoth\`ese d'extraction born\'ee. En effet, la
propri\'et\'e z\'ero-connaissance garantit $I(M;\Sigma|C) \approx 0$, de
sorte que l'hypoth\`ese d'extraction born\'ee donne
$I(V;\Sigma|C) \leq 0 + \Delta$~; seule la connaissance pr\'ealable
$\Delta$ permet alors une extraction non nulle.

\medskip
\noindent\textbf{Sch\'ema 4. La transparence active l'interpr\'etabilit\'e.}
La structure de v\'erification transparente d'ECDSA permet une
\'evaluation institutionnelle multidimensionnelle, $18$ bits extraits
sur $20$ disponibles, ce qui montre qu'une structure appropri\'ee peut \^etre
r\'ev\'el\'ee sans compromettre la solidit\'e cryptographique, tout en activant
l'utilit\'e institutionnelle.
\end{sloppypar}

% -----------------------------------------------------------------------
\subsection{Analyse de sensibilit\'e \`a la capacit\'e d'extraction}
\label{subsec:analyse_sensibilite}
% -----------------------------------------------------------------------
\begin{sloppypar}

\noindent\textbf{Impact de $\Delta$ sur l'extraction.} Le param\`etre
$\Delta$ majore l'information que l'interpr\`ete tire de la preuve
\emph{au-del\`a} de ce qu'elle dit du message. Son impact sur
$I(V;\Sigma|C)$, donc sur $\Opp$, d\'epend du r\'egime du protocole. Les
valeurs du tableau~\ref{tab:sensibilite_opp_delta} ne d\'erivent d'aucun mod\`ele
de r\'eponse. Elles illustrent des tendances plausibles, que le protocole de
la section~\ref{subsec:protocole_delta} permettrait de mesurer.

\begin{table}[htbp]
\centering
\caption{Sensibilit\'e de $\Opp$ \`a la capacit\'e d'extraction $\Delta$ (valeurs illustratives, sans mod\`ele de r\'eponse)}
\label{tab:sensibilite_opp_delta}
\small
\begin{tabular}{@{}llcccc@{}}
\toprule
\textbf{Syst\`eme} & \textbf{Sc\'enario} & $\Delta$ (bits) & $I(V;\Sigma|C)$ & $H(V|C)$ & $\Opp$ \\
\midrule
ECDSA & Standard (juge PKI) & 35 & 18 & 20 & 0,90 \\
ECDSA & Expert \'etendu & 50 & 19 & 20 & 0,95 \\
ECDSA & Novice & 15 & 10 & 20 & 0,50 \\
\midrule
Helios & Standard (juge g\'en\'eral) & 56 & 1 & 2 & 0,50 \\
Helios & Tribunal sp\'ecialis\'e & 70 & 1,4 & 2 & 0,70 \\
Helios & Novice & 30 & 0,6 & 2 & 0,30 \\
\midrule
zk-SNARK & Standard (r\'egulateur financier) & 73 & 2 & 3 & 0,67 \\
zk-SNARK & Expert cryptographe & 90 & 2,5 & 3 & 0,83 \\
zk-SNARK & R\'egulateur novice & 40 & 1,2 & 3 & 0,40 \\
\bottomrule
\end{tabular}
\end{table}

\noindent\textbf{Observation 1. Impact fort sur les syst\`emes
transparents.} Pour ECDSA, passer de $\Delta = 15$ (juge non initi\'e)
\`a $\Delta = 50$ (expert \'etendu) fait varier $\Opp$ de $0{,}50$ \`a
$0{,}95$, soit quarante-cinq points. La transparence de la structure de
preuve permet \`a l'expertise suppl\'ementaire de se traduire directement
en capacit\'e d\'ecisionnelle.

\noindent\textbf{Observation 2. Impact limit\'e sur les syst\`emes
z\'ero-connaissance.} Pour Helios, m\^eme lorsque $\Delta$ passe de
$30$ \`a $70$ bits, $\Opp$ ne varie que de $0{,}30$ \`a $0{,}70$. La
propri\'et\'e z\'ero-connaissance cr\'ee un plafond sur
l'extraction, ind\'ependamment de la capacit\'e institutionnelle, de sorte que
l'opacit\'e cryptographique ne se compense que partiellement par
l'expertise pr\'ealable. Ce plafond est compatible avec le goulot
contextuel, qui majore $I(V;\Sigma|C)$ par $H(V|C) = 2$~bits.

\noindent\textbf{Observation 3. H\'et\'erog\'en\'eit\'e des r\'egimes.}
Les trois syst\`emes appartiennent \`a des r\'egimes fondamentalement
diff\'erents selon le rapport $I(V;\Sigma|C)/H(V|C)$ avant l'apport de
$\Delta$~:
\begin{itemize}
\item ECDSA rel\`eve d'un r\'egime riche en information ($I \approx H(V|C)$ m\^eme sans $\Delta$)~;
\item Helios rel\`eve d'un r\'egime pauvre en information ($I \approx H(V|C)/2$ m\^eme avec $\Delta$ maximal)~;
\item zk-SNARK rel\`eve d'un r\'egime interm\'ediaire, sensible \`a $\Delta$ mais born\'e.
\end{itemize}
\end{sloppypar}

% -----------------------------------------------------------------------
\subsection{Implications pour la conception et la r\'eglementation}
\label{subsec:implications_conception}
% -----------------------------------------------------------------------
\begin{sloppypar}
Le Trilemme CRO conditionnel r\'eoriente la conception de protocoles
cryptographiques, qui passe de l'optimisation abstraite (\og quelles propri\'et\'es
de s\'ecurit\'e sont atteignables~? \fg) \`a l'ing\'enierie contrainte
par la capacit\'e (\og quelles propri\'et\'es cette institution peut-elle
effectivement utiliser, \'etant donn\'e sa capacit\'e mesurable
$\Delta$~? \fg).

\subsubsection{Choix de la variante SOTP}
\label{subsubsec:tableau_comparatif_sotp}\label{subsubsec:algo_selection_variante}

Les trois variantes SOTP sont compar\'ees au tableau~\ref{tab:comparaison_sotp_construction}
et l'Algorithme~\ref{alg:selection_variante_sotp} en formalise le choix
(chapitre~\ref{chap:constructions}).

L'algorithme part des deux seuils qui d\'elimitent les r\'egimes d'entropie
contextuelle, $H^*(\lambda) = \sqrt{\lambda}$ et $H^{**}(\lambda) =
\lambda/\log_2\lambda$. Il traite d'abord le cas bloquant. Si le contexte
est assez riche pour exiger SOTP-High mais que l'institution ne dispose pas
d'un trusted setup fiable, aucune recommandation optimale n'existe. La
proc\'edure signale l'impasse sans la masquer~; elle rabat alors
sur SOTP-Base, au prix d'une perte d'opposabilit\'e explicite. Dans les
autres cas, elle compare $H(\mathcal{C})$ aux deux seuils et retourne la
variante du r\'egime correspondant, soit SOTP-Low en contexte rare, SOTP-Base
en r\'egime de transition (o\`u elle dimensionne au passage la projection
al\'eatoire) ou SOTP-High en contexte riche, o\`u le circuit de preuve doit
\^etre compil\'e. Le co\^ut de la d\'ecision se r\'eduit \`a deux
comparaisons~; l'essentiel du travail reste l'estimation pr\'ealable de
$H(\mathcal{C})$ et de $\Delta$.

\subsubsection{M\'ethodologie de conception en six phases}
\label{subsubsec:methodologie_cinq_phases}

La progression F$\to$B$\to$E$\to$S$\to$C$\to$X fournit un cadre de d\'ecision explicite l\`a o\`u la
pratique actuelle consiste souvent \`a d\'eployer le protocole
cryptographiquement le plus puissant disponible, ind\'ependamment de
la capacit\'e institutionnelle de l'environnement cible.

\begin{table}[htbp]
\centering
\caption{M\'ethodologie de conception
F$\to$B$\to$E$\to$S$\to$C$\to$X}
\label{tab:methodologie_FBESCX}
\begin{tabular}{@{}p{1.8cm}p{2.8cm}p{6.5cm}@{}}
\toprule
\textbf{Phase} & \textbf{Action} & \textbf{Outils} \\
\midrule
\textbf{F} & Formaliser le contexte
  & Estimation $H(\mathcal{C})$~\cite{Shannon48,CoverThomas06} \\
\textbf{B} & Borner la capacit\'e
  & Triangulation cognitive ($\Delta$)~\cite{Miller56,Simon82,Ericsson95,Chi2006} \\
\textbf{E} & \'Etendre au r\'egime~PQ
  & \'Evaluer $\Delta_q$, $\eta_q$~(section~\ref{sec:cro_pq}) \\
\textbf{S} & Saturer~: retenir la variante SOTP du r\'egime
  & Variante SOTP (Low/Base/High), heuristique~(section~\ref{sec:sotp_construction}) \\
\textbf{C} & Certifier la faisabilit\'e
  & Exigences comparées \`a la borne raffin\'ee~\cite{MINKA2026CLO,GDPR2016,eIDAS14} \\
\textbf{X} & eXp\'erimenter la mesure
  & Protocoles exp\'erimentaux~(chapitre~\ref{chap:experimentation}) \\
\bottomrule
\end{tabular}
\end{table}

\noindent\textbf{Phase C. Test de faisabilit\'e.} V\'erifier~:
\begin{equation}
\Conf_{\min} \cdot \mathsf{Rel}_{\mathrm{rec},\min} \cdot \Opp_{\min} \stackrel{?}{\leq}
\frac{1+\kappa/\lambda}{\lambda}\cdot\min\Bigl\{1,\;\frac{\kappa/\lambda+\Delta}{\gamma}\Bigr\}
\label{eq:test_faisabilite}
\end{equation}
(forme raffin\'ee~\eqref{eq:trilemme_cro_raffine}, applicable aux protocoles qui
satisfont les hypoth\`eses du th\'eor\`eme).
Si l'in\'egalit\'e est viol\'ee, trois options s'offrent au concepteur, \`a savoir relaxer une exigence,
augmenter $\Delta$ par investissement dans la formation institutionnelle
ou augmenter $\lambda$ (impact marginal si $\Delta$ reste fixe).

\subsubsection{Strat\'egie de conception par r\'egime}
\label{subsubsec:strategie_par_regime}

\begin{definition}[R\'egime de basse capacit\'e]\label{def:regime_low}
Une institution op\`ere en r\'egime de \textbf{basse capacit\'e} si
$\Delta < 0{,}1\lambda$. Elle ne dispose pas de l'expertise n\'ecessaire
pour interpr\'eter des preuves cryptographiques
complexes~\cite{Herley09}. \textit{Recommandation.} Des m\'ecanismes
transparents (ECDSA avec pistes d'audit claires). La cryptographie
pr\'eservant la vie priv\'ee est inadapt\'ee \`a ce niveau.
\end{definition}

\begin{definition}[R\'egime de capacit\'e moyenne]\label{def:regime_mid}
Une institution op\`ere en r\'egime de \textbf{capacit\'e moyenne} si
$0{,}1\lambda \le \Delta \le 0{,}5\lambda$. Elle peut atteindre des
compromis CRO \'equilibr\'es~\cite{Grossklags10}.
\textit{Recommandation.} Des protocoles comme Helios avec des aides
\`a l'interpr\'etation structur\'ees.
\end{definition}

\begin{definition}[R\'egime de haute capacit\'e]\label{def:regime_high}
Une institution op\`ere en r\'egime de \textbf{haute capacit\'e} si
$\Delta > 0{,}5\lambda$. Elle peut soutenir des preuves cryptographiques
avanc\'ees~\cite{Grossklags10,Simon82}. \textit{Recommandation.}
Des zk-SNARKs, avec un investissement substantiel en formation et infrastructure
d'expertise.
\end{definition}

L'impl\'ementation se d\'eroule en trois \'etapes. \textbf{\'Etape~1.
Profilage institutionnel.} \'Evaluer $\Delta$ par inventaire de la
disponibilit\'e d'experts, tests de charge cognitive sur des preuves
repr\'esentatives, analyse de d\'ecisions historiques, audit de la
qualit\'e de documentation~\cite{Miller56,Ericsson95,Simon82}.
\textbf{\'Etape~2. S\'election protocolaire.} Choisir les protocoles
correspondant au r\'egime mesur\'e~; ne jamais d\'eployer un protocole
dont le r\'egime cible d\'epasse la capacit\'e actuelle sans planification
de transition. \textbf{\'Etape~3. Renforcement de capacit\'e.}
Investir dans la formation, les outils d'aide \`a la d\'ecision et les
trajectoires de migration progressives.

\subsubsection{Trois recommandations de conception}
\label{subsubsec:trois_recommandations}

L\`a o\`u la strat\'egie par r\'egime ci-dessus indique \emph{quel}
protocole retenir pour une capacit\'e donn\'ee, l'analyse de sensibilit\'e
de la sous-section~\ref{subsec:analyse_sensibilite} indique \emph{o\`u
investir} pour faire \'evoluer cette capacit\'e. Elle se traduit en trois
recommandations directement actionnables.

\noindent\textit{Recommandation 1 (ad\'equation contexte--protocole).} Le
choix du protocole doit \^etre guid\'e par le $H(\mathcal{C})$ de
l'institution cible. Un tribunal g\'en\'eral ($\Delta \approx 35$ bits) ne
tire aucun b\'en\'efice de zk-SNARK avanc\'es ($\Opp \approx 0{,}40$)
l\`a o\`u ECDSA offrirait $\Opp \approx 0{,}90$.

\noindent\textit{Recommandation 2 (investissement en formation).}
L'investissement en formation institutionnelle ($\Delta$ croissant) a des
rendements d\'ecroissants pour les syst\`emes z\'ero-connaissance mais
croissants pour les syst\`emes transparents. \`A budget de formation
\'egal, le gain d'opposabilit\'e est plus \'elev\'e sur ECDSA que sur Helios.

\noindent\textit{Recommandation 3 (compromis explicite).} Tout syst\`eme
visant $\Opp > 0{,}7$ dans un contexte \`a faible $H(\mathcal{C})$ doit
accepter une r\'eduction de $\Conf$. La fronti\`ere de Pareto
$(\Conf, \Opp)$ est born\'ee par le Trilemme CRO
conditionnel~\cite{MINKA2026CRO}, pour les protocoles qui satisfont ses
hypoth\`eses. Les valeurs chiffr\'ees de ces recommandations sont
illustratives (tableau~\ref{tab:sensibilite_opp_delta}).

\subsubsection{Seuils d'opposabilit\'e et admissibilit\'e juridique}
\label{subsubsec:seuils_opposabilite}

La m\'etrique $\Opp$ permet de proposer des seuils minimaux
diff\'erenci\'es par contexte juridictionnel~\cite{Kerr04,Casey11,Mason07}.

\begin{table}[htbp]
\centering
\caption{Seuils d'opposabilit\'e sugg\'er\'es par r\'egime juridictionnel}
\label{tab:seuils_opposabilite}
\begin{tabular}{@{}p{3.5cm}p{3.5cm}c@{}}
\toprule
\textbf{R\'egime} & \textbf{Standard de preuve} & \textbf{$\Opp_{\min}$} \\
\midrule
P\'enal & Certitude raisonn\'ee~\cite{Blackstone1769,Wigmore1983}
  & $\ge 0{,}70$ \\
Civil & Balance des probabilit\'es~\cite{Federal1988}
  & $\ge 0{,}50$ \\
Arbitrage commercial & Pr\'epond\'erance~\cite{Born2014}
  & $\ge 0{,}40$ \\
Conformit\'e r\'eglementaire & Audit interne~\cite{COSO2013}
  & $\ge 0{,}30$ \\
\bottomrule
\end{tabular}
\end{table}

Ces seuils reposent sur des estimations analytiques~; la variabilit\'e
intra-institutionnelle de $\Delta$ introduit une
incertitude~\cite{Kahneman11}~; ils constituent un guide
indicatif~\cite{Casey11}.

\noindent\textbf{Implications pour l'audit des syst\`emes de
r\'emun\'eration de l'\'Etat.} L'analyse des syst\`emes de gestion de la
solde et des pensions \`a l'\'echelle nationale, particuli\`erement dans
le contexte du Cameroun, pose un d\'efi pour l'audit
informatique en environnement institutionnel contraint. L'extraction et
l'analyse exhaustive des donn\'ees de paiement d\'epassent la v\'erification
comptable, car elles sont susceptibles de d\'evoiler
des informations sensibles pour la s\'ecurit\'e de l'\'Etat. Le cadre
CLO-CRO sugg\`ere une \emph{hypoth\`ese de travail} selon laquelle, dans les contextes o\`u
$\Delta \ll \gamma$, la forme raffin\'ee de la borne pr\'edit un plafond
d'opposabilit\'e bas pour les protocoles tr\`es confidentiels, de sorte que
l'auditeur pourrait devoir arbitrer entre la force probante de sa preuve et
la protection des donn\'ees. Cette hypoth\`ese ne d\'ecoule pas du cadre
sans mesure de $\Delta$ et de $\gamma$ dans ce contexte~; elle est \`a
\'evaluer empiriquement.

\subsubsection{Implications pour les standards NIST PQC}
\label{subsubsec:implications_nist_pqc}

L'analyse des standards NIST~PQC r\'ev\`ele des implications directes
pour le d\'eploiement~\cite{NIST.FIPS.203,NIST.FIPS.204,NIST.FIPS.205,Alagic2022}.

\begin{table}[htbp]
\centering
\caption{Implications CRO pour les standards NIST~PQC}
\label{tab:implications_nist_pqc}
\begin{tabular}{@{}p{3.2cm}p{2.5cm}p{2.8cm}p{2.8cm}@{}}
\toprule
\textbf{Standard} & \textbf{R\'egime} & \textbf{Recommandation}
  & \textbf{Pr\'ecaution} \\
\midrule
FIPS~203 (ML-KEM) & Hors champ (KEM) & --- & Pas de preuve \`a interpr\'eter~\cite{Kyber2022,Peikert2016} \\
FIPS~204 (ML-DSA) & Signature, $\eta_q = 0$ & Institutions form\'ees
  & Formation requise~\cite{Dilithium2022} \\
FIPS~205 (SLH-DSA) & Signature, $\eta_q = 0$ & Haute s\'ecurit\'e
  & Taille de signature \'elev\'ee~\cite{SPHINCS2022,Bernstein2015} \\
\bottomrule
\end{tabular}
\end{table}

Pour les institutions \`a faible capacit\'e ($\Delta_q < 30$~bits), les
standards post-quantiques n\'ecessitent des investissements significatifs
en formation et outils d'interpr\'etation avant tout d\'eploiement
effectif.

\subsubsection{Opportunit\'es de standardisation}
\label{subsubsec:opportunites_standardisation}

Trois possibilit\'es de standardisation sont identifi\'ees.
Pour \textbf{NIST~PQC}, l'int\'egration de l'opposabilit\'e comme crit\`ere
additionnel permettrait une s\'election guid\'ee par les contraintes de
d\'eploiement r\'eelles. Pour \textbf{ISO/IEC~27xxx}, une norme sur les
m\'etriques CLO fournirait un protocole standardis\'e de calcul de
$H(\mathcal{C})$ et $\Delta$. Pour la \textbf{CEMAC}, une directive sur la
cryptographie institutionnelle pour la monnaie mobile permettrait \`a six
pays d'harmoniser leurs exigences r\'eglementaires.
\end{sloppypar}

% -----------------------------------------------------------------------
\subsection{Positionnement vis-\`a-vis des travaux existants}
\label{subsec:positionnement_travaux}
% -----------------------------------------------------------------------
\begin{sloppypar}

\subsubsection{Cadres cryptographiques~: ZKP et flux d'information quantitatif}
\label{subsubsec:cadres_crypto_zkp}

Les syst\`emes zero-knowledge atteignent une efficacit\'e remarquable
mais cr\'eent une vuln\'erabilit\'e pr\'ecise dans les contextes
institutionnels, puisqu'une preuve zk-SNARK offre au r\'egulateur une
v\'erification binaire sans explication de la logique
sous-jacente~\cite{Groth16,BCTV14,PLONK20}. Le paradigme
\og ne rien r\'ev\'eler \fg{} optimise pour la confidentialit\'e mais
n\'eglige intrins\`equement l'opposabilit\'e~\cite{Bonneau16,Grimmelmann15}.
CLO formalise cette tension via $\Opp$~\cite{MINKA2026CLO}.

Le flux d'information quantitatif (QIF)~\cite{Smith09,Alvim12,Clarkson05}
mesure la fuite d'information pour la minimiser, alors que CLO cherche \`a
quantifier combien d'information est rendue interpr\'etable par
l'institution pour l'optimiser~; les deux objectifs sont sym\'etriques mais
oppos\'es. Les \'etudes cognitives calibrent $\Delta$ et
$\alpha$~\cite{Miller56,Simon82,Kahneman11}~; les pr\'edictions CLO
guident la conception des \'etudes.

\subsubsection{S\'ecurit\'e utilisable}
\label{subsubsec:securite_utilisable}

La recherche en s\'ecurit\'e utilisable a \'etabli des r\'esultats
empiriquement riches (folk models~\cite{Wash2010}, rationalit\'e de
l'ins\'ecurit\'e~\cite{Herley09}). Elle ne formalise toutefois pas
math\'ematiquement ces limites, n'en d\'erive pas de bornes d'impossibilit\'e
et ne produit pas de constructions optimales
garanties~\cite{Cranor05,Bonneau12}. La m\'etrique $\Opp$ exprime
formellement ce que les \'etudes d'utilisateurs observent
qualitativement~\cite{MINKA2026CLO}, de sorte que les deux approches se compl\`etent.

\subsubsection{Recherche juridique et responsabilit\'e algorithmique}
\label{subsubsec:recherche_juridique}

Citron et Pasquale~\cite{Citron2008,Citron14} ont argument\'e que les
syst\`emes automatis\'es doivent fournir une explicabilit\'e pour
satisfaire les exigences de due process technologique. Ces travaux
identifient pr\'ecis\'ement les probl\`emes que CLO cherche \`a
formaliser, mais s'arr\^etent \`a l'analyse
normative~\cite{Casey11,Kerr04}. CLO leur fournit un outil
quantitatif, en compl\'ement de l'analyse normative. La valeur de cette
instrumentation ne sera \'etablie qu'une fois que des \'etudes
contr\^ol\'ees auront confirm\'e que les valeurs analytiques de $\Opp$
pr\'edisent effectivement les jugements d'admissibilit\'e r\'eels.
\end{sloppypar}

% -----------------------------------------------------------------------
\subsection{Limitations et probl\`emes ouverts}
\label{subsec:limitations_problemes_ouverts}
% -----------------------------------------------------------------------
\begin{sloppypar}

\subsubsection{Limitations techniques}
\label{subsubsec:limitations_techniques}

\textbf{Les axiomes de Fano et Pinsker ne sont pas d\'eriv\'es ab initio.}
Notre formalisation les int\`egre comme propri\'et\'es des instances de
protocole sans les d\'eriver \`a partir de la th\'eorie des
probabilit\'es. Une formalisation compl\`ete n\'ecessiterait l'int\'egration
des biblioth\`eques de mesure-th\'eorie de Mathlib.

\textbf{Les constructions SOTP ne sont pas encore couvertes par Lean4.}
La formalisation couvre le noyau alg\'ebrique du Trilemme, mais pas les
r\'eductions cryptographiques des constructions SOTP ni l'extension
post-quantique. Ces extensions constituent une direction de recherche
future.

\subsubsection{Limitations m\'ethodologiques (L1--L6)}
\label{subsubsec:limitations_methodologiques}

Six limitations d\'elimitent pr\'ecis\'ement le p\'erim\`etre
d'applicabilit\'e des r\'esultats.

\paragraph{L1. Nature analytique des estimations CLO.}
Les m\'etriques $\Conf$, $\Relver$, $\Opp$ des syst\`emes
analys\'es ont \'et\'e calcul\'ees par analyse param\'etrique \`a partir
de sp\'ecifications publi\'ees, sans mesure aupr\`es d'acteurs
institutionnels r\'eels~\cite{MINKA2026CLO}. Des \'etudes contr\^ol\'ees
avec juges, auditeurs et garants en situation r\'eelle sont n\'ecessaires
pour \'etablir la correspondance entre m\'etriques analytiques et
comportement observ\'e~\cite{Casey11,Goodman15,Campbell1976,Saltelli2008}. De
plus, ces syst\`emes ne sont pas des instances du Trilemme
(Remarque~\ref{rem:portee_cas_etude}), de sorte que leurs valeurs calibrent le cadre sans le
tester.

\paragraph{L2. Quantification non standardis\'ee de $H(\mathcal{C})$.}
La m\'ethode de calcul de l'entropie contextuelle n'est pas
standardis\'ee. Deux chercheurs analysant le m\^eme syst\`eme pourraient
d\'eriver des valeurs diff\'erentes selon les hypoth\`eses retenues sur
la granularit\'e du contexte. La voie de r\'esolution passe par un
consensus communautaire.

\paragraph{L3. Statut de publication des security games.}
Les bornes CLO publi\'ees dans IEEE Access~\cite{MINKA2026CLO,MINKA2026CRO,MINKA2026PQCRO}
sont \'etablies dans le formalisme information-th\'eorique. Les jeux
cryptographiques (IND-CONF, BIND-REL, EXT-OPP) et les r\'eductions
game-based des constructions SOTP sont d\'evelopp\'es dans les
pr\'epublications~\cite{MINKA2025CRO,MINKA2025Q2CSI,MINKA2025ZKNR_UC}
mais n'ont pas encore fait l'objet d'une publication dans une revue
\`a comit\'e de lecture.

\paragraph{L4. \'Equivalence formelle incompl\`ete.}
Les ponts \'etablis entre les deux formalismes montrent que la
r\'esistance aux jeux \emph{implique} des bornes inf\'erieures sur les
m\'etriques CLO. La r\'eciproque, selon laquelle une forte opposabilit\'e
information-th\'eorique impliquerait une r\'esistance aux jeux,
n'est pas \'etablie. Cette asym\'etrie est une limitation
structurelle, puisque les garanties information-th\'eoriques sont
\emph{n\'ecessaires} mais pas \emph{suffisantes} pour la s\'ecurit\'e
computationnelle.

\paragraph{L5. Mod\`eles quantiques simplifi\'es.}
L'extension quantique utilise principalement le canal d\'epolarisant
comme mod\`ele de bruit. Des ph\'enom\`enes plus fins (amortissement
d'amplitude~\cite{Devoret2013}, d\'ephasage~\cite{Breuer2002}, effets de
m\'emoire quantique~\cite{Rivas2012}) peuvent ne pas y \^etre
captur\'es. Des mod\`eles de canaux sp\'ecifiques pourraient donner des
bornes plus fines.

\paragraph{L6. Port\'ee limit\'ee aux preuves non interactives.}
L'analyse couvre principalement les syst\`emes de preuves non
interactives~\cite{Groth16,BCTV14}. Les protocoles \`a rounds
multiples~\cite{Goldreich1987}, les architectures
multi-prouveurs~\cite{CramerDamgardSchoenmakers} et les verdicts
distribu\'es~\cite{Evans2018} peuvent pr\'esenter des caract\'eristiques
de compromis substantiellement diff\'erentes.

\subsubsection{Limitations conceptuelles}
\label{subsubsec:limitations_conceptuelles}

\textbf{Mod\`ele de contexte statique.} Nous traitons le contexte $C$
comme fixe. Or les institutions apprennent et s'adaptent, de sorte que $H(C)$ et
$\Delta$ \'evoluent par accumulation de pr\'ec\'edents, changement
technologique et mises \`a jour r\'eglementaires~\cite{Simon82}.
\textbf{Fonctions de verdict d\'eterministes.} Nous mod\'elisons
$V(\Sigma, C)$ comme d\'eterministe dans les exemples (auquel cas $\Opp = 1$,
Remarque~\ref{rem:verdict_deterministe}), alors que le th\'eor\`eme autorise un
verdict qui d\'epend aussi de ressources institutionnelles. La prise de d\'ecision r\'eelle
implique des processus stochastiques, des d\'esaccords multi-personnes et
des biais cognitifs~\cite{Casey10,Kahneman11}. \textbf{Hypoth\`ese de
$\beta$-r\'egularit\'e.} Le Th\'eor\`eme~\ref{thm:conditional_CRO_trilemma}
requiert $\beta = O(1)$~; des constructions pathologiques pourraient la
violer. \textbf{Port\'ee du t\'emoin.} Le t\'emoin de coh\'erence Lean4
atteint la borne dans le mod\`ele abstrait (forme rel\^ach\'ee de Fano), mais il
n'est pas r\'ealisable par un protocole, de sorte que la question d'un protocole concret
qui approche la borne reste ouverte.
\end{sloppypar}

\begin{boiteArbitre}
L'arbitre dresse l'inventaire de ce qui manque. Certains axiomes sont admis sans
\^etre d\'eriv\'es, tandis que la couverture Lean4 n'inclut encore ni SOTP ni
l'extension post-quantique. Six limitations m\'ethodologiques sont
nomm\'ees sans d\'etour, parmi lesquelles les estimations analytiques, l'entropie non
standardis\'ee, le r\'ealisme computationnel, le bruit quantique simplifi\'e
et la restriction aux preuves non interactives. Quatre limitations conceptuelles touchent
aussi le mod\`ele lui-m\^eme, \`a savoir le contexte statique, les verdicts d\'eterministes,
la r\'egularit\'e suppos\'ee et la port\'ee du t\'emoin.

Il note que rien n'est dissimul\'e, puisque chaque limite est pos\'ee et que chaque
cons\'equence en est tir\'ee. Il y voit un dossier en cours d'instruction,
l'instruction exigeant de savoir exactement
o\`u sont les trous.
\end{boiteArbitre}

% -----------------------------------------------------------------------
\subsection{La port\'ee constructive~: de l'impossibilit\'e au guide}
\label{subsec:portee_constructive}
% -----------------------------------------------------------------------
\begin{sloppypar}

Les th\'eor\`emes d'impossibilit\'e f\'econds sont souvent accompagn\'es de
constructions qui en montrent la port\'ee. CAP trouve son sens dans les
syst\`emes distribu\'es r\'eels (Spanner, CockroachDB)~\cite{Brewer00,GilbertLynch02}.
FLP trouve le sien dans les protocoles partiels (Paxos,
Raft)~\cite{FLP85}. Le Trilemme CRO se prolonge dans la
\textbf{famille SOTP}, qui r\'eunit trois constructions cryptographiques, une par r\'egime
d'entropie contextuelle, dont les propri\'et\'es de confidentialit\'e et de
fiabilit\'e de v\'erification sont prouv\'ees mais dont l'opposabilit\'e est
guid\'ee par une heuristique de conception.

\subsubsection{Ce que la construction apporte et ce qu'elle n'\'etablit pas}

\textbf{La borne devient utilisable pour la conception.} Un ing\'enieur qui
conna\^it la forme raffin\'ee de la borne sait quelles exigences sont
incompatibles~; avec SOTP, il dispose aussi d'un point de d\'epart par
r\'egime (contexte rare $\to$ SOTP-Low, transition $\to$ SOTP-Base, contexte
riche $\to$ SOTP-High).

\textbf{La finesse de la borne reste ouverte.} Les fractions de la borne
atteintes ($0{,}45$--$0{,}75$) rel\`event d'heuristiques et le t\'emoin de
coh\'erence Lean n'est pas r\'ealisable. Qu'un protocole concret approche la borne \`a facteur
constant est une question ouverte (Conjecture~\ref{conj:optimalite_universelle}).

\textbf{La borne devient un guide.} Elle dit
quel compromis est exclu~; les constructions indiquent comment s'en
approcher, sans garantie d'optimalit\'e.

\subsubsection{Port\'ee de la contribution SOTP}

\textbf{SOTP-Low} couvre le r\'egime des contextes rares
($H(C) < \sqrt{\lambda}$). Elle utilise un encodage d\'eterministe du contexte,
avec confidentialit\'e calculatoire sous PRG~\cite{BlumMicali84} et fiabilit\'e
de v\'erification sous DLog~\cite{DiffieHellman76}.

\textbf{SOTP-Base} couvre le r\'egime de transition. Elle utilise la projection
Johnson-Lindenstrauss~\cite{JohnsonLindenstrauss84,DasguptaGupta03}, dimensionn\'ee
pour pr\'eserver les distances entre contextes \`a $\varepsilon$ pr\`es.

\textbf{SOTP-High} couvre le r\'egime des contextes riches
($H(C) > \lambda/\log\lambda$). Elle utilise des zk-SNARKs
Groth16~\cite{Groth16,BCTV14} et un hachage Poseidon~\cite{Poseidon2019} pour
engager un contexte arbitrairement riche dans un condens\'e de 256 bits.

L'optimalit\'e universelle de la famille, qui interdirait toute am\'elioration
super-constante du produit par rapport \`a la meilleure variante, est une
conjecture (Conjecture~\ref{conj:optimalite_universelle}).

\begin{boiteArbitre}
L'arbitre comprend que la borne seule ne suffit pas, puisqu'une limite
d\'epourvue de construction ressemble \`a une porte ferm\'ee dont on ignorerait le plan. La famille SOTP
esquisse ce plan en indiquant par o\`u commencer dans chaque r\'egime, sans
pr\'etendre avoir atteint la limite. L'impossibilit\'e devient ainsi un guide.
\end{boiteArbitre}

\end{sloppypar}

% -----------------------------------------------------------------------
\subsection{Ce que la formalisation enseigne~: le\c{c}on m\'ethodologique}
\label{subsec:verification_corrige_lecon}
% -----------------------------------------------------------------------
\begin{sloppypar}

Les formulations alternatives issues de la formalisation Lean4, y compris le
cas limite qui \'etablit la n\'ecessit\'e du r\'egime de confidentialit\'e, sont
pr\'esent\'ees \`a la sous-section~\ref{subsec:lean4_erreurs}. La pr\'esente
sous-section se concentre sur la le\c{c}on m\'ethodologique qui en d\'ecoule.

\paragraph{La rigueur se v\'erifie.}
Les formulations pr\'ecis\'ees touchaient \`a la structure des preuves~: le sens
d'une in\'egalit\'e, une borne inf\'erieure tir\'ee d'une borne sup\'erieure et la
r\'esolution d'une optimisation. Ces points \'echappent facilement \`a la relecture
parce qu'ils produisent des r\'esultats \emph{plausibles}. La v\'erification
m\'ecanique a mis en \'evidence les questions de sens d'in\'egalit\'e et de
quantification. Elle ne dispense pas de l'examen des d\'efinitions, qui a conduit
\`a d'autres pr\'ecisions (distinction entre fiabilit\'e de reconstruction et
fiabilit\'e de v\'erification, d\'efinition de $\eta_q$), hors de port\'ee d'une
formalisation qui prend ces d\'efinitions pour donn\'ees.

\paragraph{La formalisation est un instrument de d\'ecouverte.}
Au-del\`a de la certification de r\'esultats d\'ej\`a prouv\'es, la formalisation
r\'ev\`ele les hypoth\`eses implicites et les points structurels \`a pr\'eciser. Ce
constat prolonge les observations de Scerri~\cite{Scerri14} et de
Wiedijk~\cite{Wiedijk2008}, selon lesquelles revisiter les mod\`eles de preuve
r\'ev\`ele des probl\`emes fondationnels que ni la th\'eorie pure ni
l'exp\'erimentation n'identifient.

\paragraph{La formalisation gagne \`a pr\'ec\'eder la publication.}
Dans notre cas, la formalisation Lean4 a \'et\'e entreprise \emph{apr\`es} la
publication. Elle a conduit \`a des formulations plus pr\'ecises qui auraient pu
figurer d\`es la soumission. Pour les r\'esultats interdisciplinaires combinant
th\'eorie de l'information, cryptographie et sciences cognitives, la
formalisation m\'ecanique a sa place dans le processus de preuve lui-m\^eme, en
amont de la publication.

\paragraph{Cons\'equences pour la pratique.}
Cette exp\'erience conduit \`a trois recommandations pour la communaut\'e.
\textbf{Pour les auteurs}, il s'agit de formaliser le noyau alg\'ebrique avant
soumission, m\^eme partiellement~; une formalisation partielle fait d\'ej\`a
appara\^itre les points structurels \`a pr\'eciser. \textbf{Pour les revues}, il
convient d'inclure syst\'ematiquement un relecteur \og v\'erification formelle \fg{}
pour les articles dont le noyau math\'ematique est formalisable~; la charge est
faible, le b\'en\'efice \'elev\'e. \textbf{Pour la formation}, il importe d'enseigner
la formalisation comme partie int\'egrante du m\'etier de chercheur. Il s'agit
d'une comp\'etence transf\'erable qui b\'en\'eficie \`a toutes les disciplines
math\'ematis\'ees.

\begin{boiteArbitre}
L'arbitre referme le dossier. Des formulations plus pr\'ecises sont issues de la
machine ou d'une relecture attentive des d\'efinitions, puis retenues par les
auteurs. Il en retient que ni la relecture seule ni la machine seule ne
suffisent, car l'une produit des preuves plausibles tandis que l'autre certifie
ce qu'on lui donne. L'arbitre comprend que la formalisation est une n\'ecessit\'e
\'epist\'emologique pour tout r\'esultat interdisciplinaire dont aucune des
disciplines constituantes ne contr\^ole l'int\'egralit\'e.
\end{boiteArbitre}

\end{sloppypar}

% -----------------------------------------------------------------------
\subsection{Le statut des protocoles exp\'erimentaux~: un programme
falsifiable}
\label{subsec:statut_protocoles_experimentaux}
% -----------------------------------------------------------------------
\begin{sloppypar}

Le chapitre~\ref{chap:experimentation} occupe une position
inhabituelle dans l'\'economie de cette th\`ese, puisqu'il sp\'ecifie des
protocoles exp\'erimentaux sans les avoir d\'eploy\'es. Cette position
m\'erite d'\^etre clarifi\'ee, car elle d\'etermine la port\'ee
\'epist\'emique de l'ensemble du travail.

\subsubsection{Statut \'epist\'emique}

Le chapitre~\ref{chap:experimentation} se situe entre r\'esultats
empiriques (aucun protocole n'a \'et\'e d\'eploy\'e) et th\'eorie pure (les
protocoles sont sp\'ecifi\'es jusqu'\`a l'op\'erationnel). Il rel\`eve de la
conception m\'ethodologique, puisqu'il sp\'ecifie \emph{comment} valider empiriquement
le cadre, sans encore l'avoir fait. Cette sp\'ecification constitue une
contribution en soi, pour trois raisons.

\textbf{Elle rend le cadre falsifiable.} Un cadre th\'eorique
conditionnel \`a des param\`etres non mesurables est non falsifiable au
sens de Popper~\cite{Popper1934}. En sp\'ecifiant les protocoles de
mesure, les seuils de falsification et les pr\'edictions testables, le
chapitre~\ref{chap:experimentation} transforme le Trilemme CRO en
programme scientifique falsifiable. Cette transformation conditionne sa
port\'ee empirique.

\textbf{Elle guide la communaut\'e.} Les protocoles exp\'erimentaux sont
sp\'ecifi\'es \`a un niveau op\'erationnel rare dans la litt\'erature
th\'eorique, qu'il s'agisse de la stratification des \'echantillons, des estimateurs avec biais
et variance, des crit\`eres de fiabilit\'e inter-\'evaluateurs ou du calendrier de
d\'eploiement. Cette sp\'ecification permet \`a d'autres chercheurs de
reproduire, de critiquer ou d'\'etendre le programme.

\textbf{Elle documente les limitations.} Le volet~B du chapitre d\'erive
les valeurs d'opposabilit\'e des trois syst\`emes de r\'ef\'erence et
\'etablit la coh\'erence interne du cadre. Il fournit une premi\`ere
\textbf{calibration analytique} en attente de validation
exp\'erimentale tout en documentant explicitement ce qui reste \`a mesurer.

\subsubsection{Parall\`eles avec d'autres disciplines}

Ce statut se retrouve dans les disciplines o\`u la th\'eorie pr\'ec\`ede
l'exp\'erience.

\textbf{Physique th\'eorique.} Une th\'eorie qui pr\'edit des faits
falsifiables et sp\'ecifie comment les tester est une contribution,
m\^eme si les tests n'ont pas encore \'et\'e men\'es. La pr\'ediction du
boson de Higgs a pr\'ec\'ed\'e sa d\'etection de plusieurs d\'ecennies.

\textbf{Biologie th\'eorique.} Les mod\`eles de plissement
prot\'eique~\cite{Anfinsen1973} ont pr\'ec\'ed\'e les m\'ethodes de
validation exp\'erimentale (AlphaFold~\cite{Jumper2021}).

\textbf{\'Economie comportementale.} La th\'eorie des
perspectives~\cite{Kahneman1979} a sp\'ecifi\'e des protocoles
exp\'erimentaux avant leur d\'eploiement.

Dans tous ces cas, la \textbf{sp\'ecification du programme de
validation} est reconnue comme une contribution scientifique. Il
en va de m\^eme pour le chapitre~\ref{chap:experimentation}.

\subsubsection{Cons\'equences pour la lecture de la th\`ese}

Cette clarification engage trois mani\`eres de lire ce qui pr\'ec\`ede.

\textbf{Les bornes sont conditionnelles et leurs conditions sont
\'enonc\'ees.} Le Trilemme CRO, classique comme post-quantique, ne dit rien
en dehors du r\'egime o\`u il est \'etabli, lequel suppose l'extraction born\'ee, un verdict
efficace, la r\'egularit\'e du protocole et le r\'egime de fuite
$I(M;\Sigma|C) \leq \kappa/\lambda$. Ces conditions d\'elimitent le domaine de validit\'e du
r\'esultat, de sorte qu'un protocole plac\'e hors de ce domaine \'echappe \`a la
borne. Lire le trilemme comme une interdiction absolue serait donc le
mal lire. Sa validit\'e math\'ematique est \'etablie~; sa port\'ee
empirique d\'epend de grandeurs qui restent \`a mesurer.

\textbf{Le mod\`ele est r\'efutable et les seuils de r\'efutation sont
publics} (tableau~\ref{tab:criteres_refutation}). Si les mesures de $\Delta$
ou de $\Delta_q$ s'\'ecartent syst\'ematiquement des plages attendues, ou si le
produit observ\'e d\'epasse la borne sur un syst\`eme qui satisfait les
hypoth\`eses du th\'eor\`eme, le mod\`ele
d'extraction born\'ee devra \^etre r\'evis\'e, sans r\'einterpr\'etation des
mesures. Le chapitre~\ref{chap:experimentation} fixe ces seuils
avant toute campagne de mesure, ce qui interdit de les ajuster apr\`es
coup.

\textbf{La coh\'erence \'etablie est interne.} Les ordres
de grandeur retenus pour les syst\`emes de r\'ef\'erence sont mutuellement
coh\'erents et compatibles avec la litt\'erature cognitive~; les deux voies
d'estimation ne sont pas ind\'ependantes. Cela ne constitue pas une
validation exp\'erimentale, que cette th\`ese ne revendique d'ailleurs pas.

\textbf{La port\'ee du travail est donc quadruple.} Elle est
th\'eorique par les m\'etriques et les bornes (C1--C3),
constructive par les protocoles qui les illustrent (C4),
certifi\'ee par la v\'erification m\'ecanique de leur noyau (C5) et
exp\'erimentale par le programme falsifiable qui les rend
testables (C6). Cette quadruple port\'ee structure la conclusion
g\'en\'erale.

\begin{boiteArbitre}
L'arbitre comprend que sp\'ecifier comment tester une th\'eorie est
aussi important que la th\'eorie elle-m\^eme. Une th\'eorie non testable
est une sp\'eculation. Une th\'eorie testable, m\^eme non encore
test\'ee, est une science. Le
chapitre~\ref{chap:experimentation} est ce qui fait passer le
Trilemme CRO de la premi\`ere cat\'egorie \`a la seconde.
\end{boiteArbitre}

\end{sloppypar}

% -----------------------------------------------------------------------
\subsection{Le\c{c}ons m\'ethodologiques}
\label{subsec:lecons_methodologiques}
% -----------------------------------------------------------------------
\begin{sloppypar}
Les difficult\'es rencontr\'ees produisent des enseignements
transf\'erables \`a la communaut\'e travaillant \`a l'intersection de
la cryptographie, de la th\'eorie de l'information et des m\'ethodes
formelles.

\subsubsection{La formalisation renforce les th\'eor\`emes interdisciplinaires}
\label{subsubsec:formalisation_renforce}

La v\'erification formelle expose les hypoth\`eses cach\'ees et clarifie
les d\'ependances logiques, l\`a o\`u le raisonnement informel les
absorbe tacitement~\cite{Wiedijk2008,Avigad2014}. Ce r\'esultat est
coh\'erent avec Scerri~\cite{Scerri14}, pour qui revisiter les mod\`eles de
preuve r\'ev\`ele des probl\`emes fondationnels que ni la th\'eorie pure
ni l'exp\'erimentation n'identifient.

\subsubsection{L'axiomatisation est une n\'ecessit\'e pragmatique}
\label{subsubsec:axiomatisation_pragmatique}

Tant que les math\'ematiques sous-jacentes ne sont pas enti\`erement
formalis\'ees dans une biblioth\`eque partag\'ee, int\'egrer les
r\'esultats profonds comme propri\'et\'es d'instance est une approche
valide et transparente. La cl\'e est de s\'eparer nettement le noyau
certifi\'e des fondations accept\'ees comme
vraies~\cite{BartheEA2013,Doussot23}.

\subsubsection{Les constantes explicites sont n\'ecessaires}
\label{subsubsec:constantes_explicites}

Masquer les constantes dans $O(\cdot)$ est un raccourci acceptable en
s\'ecurit\'e classique~; en cryptographie institutionnelle, o\`u
$\Delta$, $\alpha$, $\beta$, $H(C)$ doivent \^etre mesur\'es pour
permettre la d\'ecision juridique et administrative, cette abstraction
devient un obstacle op\'erationnel. Notre borne explicite
$\gamma^{-1}(\Delta/\lambda + \kappa(\Delta+1)/\lambda^2 +
\kappa^2/\lambda^3)$ rend calculable la borne sur le produit CRO pour
tout ensemble de param\`etres institutionnels mesur\'es.

\subsubsection{La transparence scientifique comme mod\`ele}
\label{subsubsec:transparence_scientifique}

Judson~\cite{Judson22} note un sch\'ema pr\'eoccupant, \`a savoir que les mises en garde
et les d\'ecisions de conception difficiles sont souvent renvoy\'ees aux
praticiens \`a travers des heuristiques mal comprises. Notre
documentation explicite des bornes non prouv\'ees, des hypoth\`eses
n\'ecessaires et des limitations empiriques r\'epond \`a cette
pr\'eoccupation, car la rigueur consiste \`a d\'elimiter pr\'ecis\'ement les probl\`emes
ouverts au lieu de les \'eviter.

\subsubsection{La v\'erification formelle comme n\'ecessit\'e \'epist\'emologique}
\label{subsubsec:verification_necessite}

La formalisation Lean4 constitue une n\'ecessit\'e
\'epist\'emologique pour les r\'esultats interdisciplinaires combinant
th\'eorie de l'information, sciences cognitives et cryptographie, dont
aucun des trois champs n'a la ma\^itrise compl\`ete. En for\c{c}ant chaque
hypoth\`ese, quantification et approximation \`a \^etre explicite, la
v\'erification m\'ecanique approfondit la compr\'ehension du
th\'eor\`eme lui-m\^eme. Les formulations alternatives et le cas limite sont
document\'es \`a la sous-section~\ref{subsec:lean4_erreurs}.
\end{sloppypar}

\begin{boiteArbitre}
L'arbitre retient quatre certitudes. La cryptographie institutionnelle
constitue un domaine propre, car sa question centrale n'\'etait pos\'ee nulle part.
Une impossibilit\'e prouv\'ee repr\'esente la forme la plus nette de
connaissance. Ce qu'il ne peut pas obtenir simultan\'ement a
d\'esormais un nom, une mesure et une borne. Enfin, la rigueur formelle a
mis au jour des propri\'et\'es
informationnelles transform\'ees en conditions v\'erifiables, une
structure de preuve pr\'eservant la borne dominante et des constantes que
la notation $O(\cdot)$ dissimulait. Il sait maintenant qu'il attend
de la section suivante les perspectives qui transforment ces contraintes
en programme.
\end{boiteArbitre}

\begin{table}[htbp]
\centering
\caption{Synth\`ese~: r\'ev\'elations, implications, limitations et
le\c{c}ons par contribution}
\label{tab:synthese_discussions}
\begin{tabular}{@{}p{1.8cm}p{3cm}p{3cm}p{3.5cm}@{}}
\toprule
\textbf{Contrib.} & \textbf{R\'ev\'elation / R\'esultat}
  & \textbf{Implication pratique} & \textbf{Travail futur} \\
\midrule
C1~: CLO & M\'etrique $\Opp$ in\'edite~\cite{MINKA2026CLO}
  & Quantification des compromis
  & \'Etudes contr\^ol\'ees \\
C2~: Trilemme & Borne $\gamma^{-1}(\Delta/\lambda + \kappa(\Delta+1)/\lambda^2 + \kappa^2/\lambda^3)$~\cite{MINKA2026CRO}
  & Guide conception classique
  & Extension Lean4 \\
C3~: PQ & Borne $\gamma_q^{-1}(\Delta_q/\lambda + \kappa(\Delta_q+1)/\lambda^2 + \kappa^2/\lambda^3)$~\cite{MINKA2026PQCRO};
  d\'ecomp. $\delta_{\text{enh}},\delta_{\text{degr}}$
  & Standards NIST~PQC
  & Canaux quantiques r\'eels \\
C4~: SOTP & $\Conf$, $\Relver$ prouv\'ees~; opposabilit\'e heuristique~(section~\ref{sec:sotp_construction})
  & Protocoles de r\'ef\'erence
  & Preuves interactives \\
C5~: Lean4 & Propri\'et\'es d'instance explicites~;
  forme raffin\'ee certifi\'ee~(section~\ref{sec:lean4})~;
  t\'emoins (mod\`ele abstrait)
  & Certification m\'ecanique
  & Extension SOTP et PQ \\
C6~: Protocoles exp\'erimentaux & Programme de validation falsifiable
  & Feuille de route empirique
  & D\'eploiement effectif \\
\midrule
\multicolumn{4}{l}{\textit{Limitations transversales~:}} \\
L1--L6 & Estimations analytiques~; $H(\mathcal{C})$ non standardis\'e~;
  r\'ealisme computationnel~; bruit quantique simplifi\'e~;
  preuves non interactives
  & \multicolumn{2}{l}{Programme exp\'erimental
    futur} \\
\bottomrule
\end{tabular}
\end{table}